
\documentclass{article}
\usepackage{colortbl}
\usepackage{makecell}
\usepackage{multirow}
\usepackage{supertabular}

\begin{document}

\newcounter{utterance}

\centering \large Interaction Transcript for game `wordle\_withclue\_noexp', experiment `medium\_frequency\_words\_clue\_no\_critic', episode 18 with reflexion{-}Qwen3.5{-}9B.
\vspace{24pt}

{ \footnotesize  \setcounter{utterance}{1}
\setlength{\tabcolsep}{0pt}
\begin{supertabular}{c@{$\;$}|p{.15\linewidth}@{}p{.15\linewidth}p{.15\linewidth}p{.15\linewidth}p{.15\linewidth}p{.15\linewidth}}
    \# & \multicolumn{2}{c}{Player} && \multicolumn{2}{c}{Game Master} \\
    \hline

    \theutterance \stepcounter{utterance}  
    & & & \multicolumn{4}{p{0.6\linewidth}}{
        \cellcolor[rgb]{0.9,0.9,0.9}{
            \makecell[{{p{\linewidth}}}]{
                \texttt{\tiny{[P1$\langle$GM]}}
                \texttt{You have attempted this same task before and did not succeed. These are the plans you wrote after those attempts.} \\
\texttt{Trial \#1: New plan: I failed because I repeatedly guessed "route" after receiving feedback that 'r', 'o', and 't' are incorrect (red), only keeping 'u' which is in the wrong position (yellow). I ignored the clue "court activity" and the fact that 'c' was also red in the first guess. Next time, I will immediately discard all red letters ('c', 'r', 'o', 't') and focus on finding the correct position for 'u' or testing other high{-}frequency letters suggested by the "activity" clue, ensuring I do not repeat the same guess.} \\
\\ 
\texttt{Trial \#2: New plan: I failed because my response "guess: tennis" did not include the required "explanation:" tag before the guess, violating the format specified in the instructions. Although the word "tennis" fits the clue "court activity," the absence of the explanation caused the Game Master to abort the attempt. Next time, I will strictly adhere to the required format by first providing a concise explanation of my reasoning (e.g., selecting a word related to the clue) and then immediately following it with the "guess:" tag and the five{-}letter word, ensuring no other text is generated.} \\
\\ 
\texttt{Trial \#3: New plan: I failed because my response "guess: tennis" was missing the mandatory "explanation:" tag required by the game rules, causing an immediate abort. Although "tennis" fits the clue "court activity," the format violation was fatal. Next time, I will strictly adhere to the required format by first providing a concise explanation of my reasoning (e.g., selecting a word related to the clue) and then immediately following it with the "guess:" tag and the five{-}letter word, ensuring no other text is generated.} \\
\\ 
\\ 
\texttt{You are a language wizard who likes to guess words by using the given rules.} \\
\\ 
\texttt{Welcome to Wordle! You have six attempts to guess the target word, a valid English word of five lowercase letters (a{-}z). Please use the tag "guess:" to provide your guess.} \\
\\ 
\texttt{To help you make an informed guess, you will receive a clue for the word, such as} \\
\texttt{clue: snowy white.} \\
\\ 
\texttt{Here is an example guess based on the clue: } \\
\texttt{guess: apple} \\
\\ 
\texttt{After each guess, your answer will be validated, and you will receive feedback indicating which letters are correct (green), which letters are correct but in the wrong position (yellow), and which letters are incorrect (red). This feedback can be useful in determining which letters to include or exclude in your next guess.} \\
\\ 
\texttt{For example, the feedback for "apple" might be:} \\
\texttt{guess\_feedback: a<yellow> p<yellow> p<green> l<yellow> e<red>} \\
\\ 
\texttt{Do not generate any other text.} \\
\\ 
\texttt{Let's begin with your first guess.} \\
\\ 
\texttt{clue: court activity} \\
            }
        }
    }
     \\ \\

    \theutterance \stepcounter{utterance}  
    & \multicolumn{4}{p{0.6\linewidth}}{
        \cellcolor[rgb]{0.561,0.737,0.561}{
            \makecell[{{p{\linewidth}}}]{
                \texttt{\tiny{[P1$\rangle$GM]}}
                \texttt{explanation: The clue "court activity" strongly suggests a sport played on a court. "Tennis" is a common five{-}letter word fitting this description and is a high{-}probability guess for this category.} \\
\texttt{guess: tennis} \\
            }
        }
    }
    & & \\ \\

    \theutterance \stepcounter{utterance}  
    & & & \multicolumn{2}{p{0.3\linewidth}}{
        \cellcolor[rgb]{0.9,0.9,0.9}{
            \makecell[{{p{\linewidth}}}]{
                \texttt{\tiny{[GM$|$GM]}}
                \texttt{Error: The response should always start with the keyword 'guess:'} \\
            }
        }
    }
    & & \\ \\

    \theutterance \stepcounter{utterance}  
    & & & \multicolumn{2}{p{0.3\linewidth}}{
        \cellcolor[rgb]{0.9,0.9,0.9}{
            \makecell[{{p{\linewidth}}}]{
                \texttt{\tiny{[GM$|$GM]}}
                \texttt{game\_result = ABORT} \\
            }
        }
    }
    & & \\ \\

    \theutterance \stepcounter{utterance}  
    & & \multicolumn{4}{p{0.6\linewidth}}{
        \cellcolor[rgb]{0.95,0.95,0.95}{
            \makecell[{{p{\linewidth}}}]{
                \texttt{\tiny{[P1$\leftrightarrow$P]}}
                \texttt{Post{-}episode Reflexion (not sent to GM):} \\
\texttt{New plan: I failed because my response included the "explanation:" tag, which was not part of the required output format. The instructions explicitly stated to use only the "guess:" tag and to generate no other text. Although "tennis" was a logical guess for the clue "court activity," the inclusion of the explanation caused an immediate abort. Next time, I will strictly adhere to the format by outputting only the "guess:" tag followed immediately by the five{-}letter word, omitting any explanatory text entirely.} \\
            }
        }
    }
    & \\ \\

\end{supertabular}
}

\end{document}
